\section{Sensitivity and Safety Limits}
\label{sec:p3}
The propellant estimate depends on quantities not fully specified by the problem. We examine effective engine efficiency and crossflow drag while retaining the reference initial direction and 5 t planned reserve. Each known-parameter design selects a convenient recovery site, as permitted by the problem, and holds that site fixed during replay. These separate designs do not establish recovery from unknown parameter errors at one prepositioned ship. The larger influence of the unspecified separation direction is reported in Section~\ref{sec:separation-direction}.

\subsection{Dependence of the Propellant Budget on Assumptions}
\subsubsection{Effective Specific Impulse}
We compare effective specific impulses of 300, 330, and 360 s, with the thrust, throttle range, gimbal limit, aerodynamic coefficients, and mixture ratio unchanged. Each case starts from the spherical reference and reoptimizes the inventory, burn timing, controls, and ship position for a single landing burn. All three solutions meet the touchdown limits in replays at 0.005 and 0.0025 s; Table~\ref{tab:p3-isp-budget} lists the finer replay.

\begin{table}[htbp]
\centering
\caption{Effective-$I_{sp}$ comparison on the spherical model. Each design reoptimizes the inventory, controls, and ship position. Masses in tonnes, rounded independently. The last column is the earlier planar result.}
\label{tab:p3-isp-budget}
\begin{tabular}{@{}rrrrrrr@{}}
\toprule
$I_{sp}$ (s) & Inventory & Consumed & Retained & Methane & Oxygen & Planar \\
\midrule
300 & 28.48 & 23.41 & 5.08 & 6.33 & 22.15 & 28.78 \\
330 & 25.97 & 20.89 & 5.08 & 5.77 & 20.20 & 26.23 \\
360 & 23.98 & 18.90 & 5.07 & 5.33 & 18.65 & 24.20 \\
\bottomrule
\end{tabular}
\end{table}

\begin{table}[htbp]
\centering
\caption{Crossflow-coefficient comparison on the spherical model, axial coefficient fixed at 0.8. The last column is the earlier planar result; a dash marks a case without a planar solution.}
\label{tab:p3-drag}
\begin{tabular}{@{}rrrrr@{}}
\toprule
$C_n$ & Inventory (t) & Peak $q$ (kPa) & Peak specific force ($g_0$) & Planar (t) \\
\midrule
0.8 & 25.94 & 295.5 & 8.70 & -- \\
1.0 & 25.96 & 293.7 & 8.21 & -- \\
1.2 & 25.97 & 292.9 & 7.84 & 26.23 \\
1.6 & 25.49 & 290.2 & 7.32 & 27.13 \\
\bottomrule
\end{tabular}
\end{table}


Relative to 330 s, the inventory rises by 2.52 t (9.7\%) at 300 s and falls by 1.99 t (7.7\%) at 360 s. Scaling the burned mass inversely with specific impulse would give 28.06 and 24.23 t instead. The difference arises because the inventory also carries the fixed 5 t reserve and because a heavier booster descends differently, so the whole trajectory has to be replanned. The planar values in the last column follow the same trend and lie 0.2--0.3 t higher. Holding thrust and mixture ratio fixed while changing specific impulse is a simplification; a real engine with different efficiency would differ in other respects too.

\subsubsection{Crossflow Drag}
Changing the crossflow coefficient affects aerodynamic braking, pitching moment, and rotational damping together. With the axial coefficient fixed at 0.8, we compare crossflow coefficients of 0.8, 1.2, and 1.6 using the same local optimization method and touchdown requirements.

\begin{table}[htbp]
\centering
\caption{Crossflow comparison. Inventories are assigned only to accepted candidates. The unresolved row is a search limitation, not proof that landing is physically impossible.}
\label{tab:p3-drag}
\begin{tabular}{@{}llrr@{}}
\toprule
$C_n$ & Outcome & Inventory (t) & Retained (t) \\
\midrule
0.8 & No accepted candidate & -- & -- \\
1.2 & Nominal reference & 26.228 & 5.078 \\
1.6 & Accepted local candidate & 27.130 & 5.072 \\
\bottomrule
\end{tabular}
\end{table}

At $C_n=1.6$, inventory increases by 0.901 t (3.44\%), while peak specific force decreases from approximately $7.703\,g_0$ to $7.187\,g_0$. This local result trades greater inventory for a lower peak load; larger crossflow drag does not automatically imply lower fuel demand. At $C_n=0.8$, the local solve reaches its 140 s limit without an accepted trajectory. We assign no fuel requirement to its failed candidate and do not infer a feasible propellant range across the full coefficient interval.


\subsection{Loads Along the Nominal Descent}
\label{sec:flight-exposure}
A safe descent depends on the loads during the whole flight, not only at touchdown. Table~\ref{tab:flight-exposure} uses the same 25.97 t spherical trajectory as the propellant answer, so loads and budget refer to one design. Dynamic pressure is $q=\rho U^2/2$, where $U$ is the speed relative to the air~\cite{dynamicPressure}. The specific force $\lVert\mathbf F_a+\mathbf T\rVert/(m g_0)$ is the non-gravitational acceleration that the structure actually carries.

\begin{table}[htbp]
\centering
\caption{Peak loads on the spherical reference trajectory, from the replay at a 0.0025 s step. Height is the radial altitude of the center of mass. The planar design gave 297.3 kPa, $7.70\,g_0$, and 9.62 MN.}
\label{tab:flight-exposure}
\begin{tabular}{@{}lrrl@{}}
\toprule
Quantity & Time (s) & Height (km) & Peak \\
\midrule
Dynamic pressure & 175.725 & 9.834 & 292.9 kPa \\
Specific-force magnitude & 187.413 & 3.343 & $7.84\,g_0$ \\
Normal aerodynamic force & 187.523 & 3.302 & 9.98 MN \\
\bottomrule
\end{tabular}
\end{table}

Dynamic pressure peaks before the engine lights at 181.32 s, while the specific force peaks during the burn. Near that second peak the body axis is about $20^\circ$ from the direction of the oncoming air, which produces a large side force. The attitude therefore matters as much as the speed and the air density. Because structural and thermal capacity are not modeled, the table states what a booster flying this trajectory would have to withstand; it does not show that it can.

\subsection{Aerodynamic Heating Estimate}
\label{sec:heating}
To gauge the heating, we apply the Sutton--Graves correlation for the convective heat flux at a stagnation point in air~\cite{suttonGraves},
\begin{equation}\label{eq:sutton}
 \dot q_s=k\sqrt{\frac{\rho}{R_n}}\,U^3,\qquad k=1.7415\times10^{-4}\ \mathrm{kg^{1/2}\,m^{-1}},
\end{equation}
along the reference trajectory, where $R_n$ is the effective nose radius of the windward surface. The booster changes attitude during the descent and meets the air partly tail first, so $R_n$ is uncertain. We bracket it between 1 m, for a rounded edge, and the 3.5 m body radius. The peak occurs at 168.7 s, at 15.5 km altitude and 1500 m/s, before the dynamic-pressure peak. It is 249 kW/m$^2$ for $R_n=1$ m and 133 kW/m$^2$ for $R_n=3.5$ m, and the heat absorbed per unit area over the whole flight is 8.6 and 4.6 MJ/m$^2$. Equation~\eqref{eq:sutton} is intended for blunt bodies at hypersonic speed, and the separation speed is only about Mach 5, so these values indicate the order of magnitude. They are the design heating that a thermal protection system would have to handle; sizing that system is outside this study.

\subsection{Cost of Limiting Dynamic Pressure}
\label{sec:qcap}
The 293 kPa peak occurs while the engine is off, so a later correction cannot reduce it. To see what a lower structural limit would cost, we add the constraint $q\le q_{max}$ to the spherical optimization at every mesh point, with a 1\% margin. Starting trajectories come from the feedback search with one extra braking burn high in the descent, followed by a second coast and the landing burn. A solution is accepted only if both replays meet every touchdown limit and the replayed peak stays at or below $q_{max}$. We have no sourced load limit for this vehicle, so the 200 and 150 kPa caps are examples.

\begin{table}[htbp]
\centering
\caption{Dynamic-pressure caps for the $15^\circ$ spherical case. Each row is the best accepted solution for that limit. Masses in tonnes.}
\label{tab:qcap}
\begin{tabular}{@{}lllrrrr@{}}
\toprule
Cap (kPa) & Status & Burns & Inventory & Entry burn & Landing burn & Peak $q$ (kPa) \\
\midrule
none & optimized & 1 & 25.97 & -- & 20.89 & 292.9 \\
200 & optimized & 2 & 37.30 & 13.50 & 18.72 & 198.0 \\
150 & optimized & 2 & 50.17 & 27.08 & 18.05 & 148.7 \\
\bottomrule
\end{tabular}
\end{table}

A 200 kPa cap raises the separation inventory from 25.97 t to 37.30 t. The added propellant goes into a 38 s entry burn starting at 43.2 km altitude, while the landing burn still uses about 19 t. A 150 kPa cap needs 50.17 t, with a 77 s entry burn starting at 62.7 km altitude. Lowering the limit from 200 to 150 kPa costs another 12.9 t. These solutions pass the same replay checks as the reference. They were still improving when the optimization was stopped, so they are upper estimates of the minimum.


At the masses involved, the 2.846 MN engine decelerates the booster at roughly $1.5$--$2\,g_0$, far less than aerodynamic drag near peak dynamic pressure. Lowering $q$ therefore takes a long burn while the booster is still fast, and all the propellant burned early has to be carried from separation. If a structural limit near or below 200 kPa applied, it rather than the landing burn would set the propellant requirement.


\subsection{A Limited Guidance Extension}
Holding the nominal commands fixed does not generally preserve touchdown under wind or initial-attitude errors. An illustrative 30 t feedback design recovered five tested steady-wind errors, $-0.1$, $-0.05$, 0, $+0.05$, and $+0.1$ m/s, at the same ship under exact state information and instantaneous actuator response. Its nominal peak specific force was approximately $8.13\,g_0$, above the reference value. This supplementary calculation demonstrates local corrective action, not a general recovery envelope. The trial inventory is not a replacement for the nominal propellant estimate. Detailed replanning, wind-error sweeps, and controller settings are retained in the repository's supplementary calculation records.

\subsection{Conclusions of the Sensitivity Analysis}
\label{sec:p3-conclusions}
The separation inventory is conditional on the selected propulsion, aerodynamics, recovery plan, and retained propellant. For the reference design, the P2 candidate carries approximately 26.23 t, including the 5 t planned usable reserve, with approximately 5.83 t methane and 20.40 t oxygen. The known-efficiency designs give accepted local inventories of 24.20--28.78 t. The higher-crossflow design requires approximately 27.13 t, while the lower-crossflow search remains unresolved. These individual designs use separately selected ship positions and do not define a common inventory or one command history valid across uncertain conditions.

At a fixed ship, the original nominal commands are sensitive to initial and wind errors. Case-specific separation-time planning restores the eleven examined cases at a common 30 t trial inventory when the initial state and steady wind are accurately known. With wind unknown to the planner, the 20 ms feedback scheme restores the five examined points from $-0.1$ to $+0.1$ m/s and retains more than 6.03 t. This sampled result supports local correction under exact state information and ideal actuator response; it does not establish a continuous wind envelope or robustness to combined errors.

Touchdown recovery and low flight loads are separate requirements. The nominal feedback trajectory reaches approximately 304 kPa and $8.13\,g_0$, above the corresponding P2 peaks. Pre-ignition pressure and flow-energy peaks cannot be reduced by changing only the later feedback. Structural and thermal safety remain undetermined without vehicle allowables and material-response models. The present results therefore provide conditional propellant budgets and specific limits of the recovery calculations. Additional thrust, reserve, deck-motion, actuator, and three-dimensional studies are possible extensions, rather than prerequisites for the bounded sensitivity analysis reported here.

